%%%%%%%%%%%%%%%%%%%%%%%%%%%%%%%%%
%
%       EXERCÍCIO
%
%%%%%%%%%%%%%%%%%%%%%%%%%%%%%%%%%

\definevalorquestao

\begin{exercicioBanco}[\valorquestao]
    O tempo de duração, em horas, de uma lâmpada especial foi modelado por uma variável aleatória \(X\) com a seguinte função de probabilidade:
     \[
        \begin{array}{c|c|c|c|c|c|c}
            X & 5 & 6 & 7 & 8 & 9 & 10 \\
            \hline
            p_{i} & 0,1 & 0,1 & 0,2 & 0,4 & 0,1 & 0,1
        \end{array}
    \]
    Cada lâmpada custa ao fabricante R\$10,00, mas se sua duração for inferior a \(6\) horas, ele se compromete a indenizar o comprador com R\$15,00. Qual deve ser o preço de cada lâmpada para o fabricante obter um lucro médio por lâmpada de R\$20,00?
\end{exercicioBanco}

